%% jorge_pizarro_cv_ats.tex
%% Compilar con: pdflatex jorge_pizarro_cv_ats.tex
%% Requiere paquetes: geometry, titlesec, enumitem, hyperref, fontenc, inputenc

\documentclass[a4paper,11pt]{article}

% ── Codificación y fuente ──────────────────────────────────────
\usepackage[T1]{fontenc}
\usepackage[utf8]{inputenc}
\usepackage{lmodern}

% ── Márgenes ──────────────────────────────────────────────────
\usepackage[top=1.8cm, bottom=1.8cm, left=1.8cm, right=1.8cm]{geometry}

% ── Colores ───────────────────────────────────────────────────
\usepackage[dvipsnames]{xcolor}
\definecolor{azulcv}{HTML}{1F4E79}
\definecolor{gris}{HTML}{555555}
\definecolor{grisCLaro}{HTML}{AAAAAA}

% ── Tipografía de sección ──────────────────────────────────────
\usepackage{titlesec}
\titleformat{\section}
  {\color{azulcv}\normalfont\bfseries\large}
  {}{0em}
  {\MakeUppercase}
  [{\color{azulcv}\vspace{2pt}\hrule height 0.8pt\vspace{4pt}}]
\titlespacing*{\section}{0pt}{10pt}{4pt}

% ── Listas ────────────────────────────────────────────────────
\usepackage{enumitem}
\setlist[itemize]{
  leftmargin=1.4em,
  topsep=2pt,
  itemsep=2pt,
  parsep=0pt,
  label={\textbullet}
}

% ── Hiperlinks ────────────────────────────────────────────────
\usepackage[hidelinks]{hyperref}

% ── Sin sangría, espaciado entre párrafos ──────────────────────
\setlength{\parindent}{0pt}
\setlength{\parskip}{3pt}

% ── Tabla para alinear fechas a la derecha ────────────────────
\usepackage{tabularx}
\usepackage{array}

% ── Sin numeración de página ──────────────────────────────────
\pagestyle{empty}

% ─────────────────────────────────────────────────────────────
% Comandos personalizados
% ─────────────────────────────────────────────────────────────

% Entrada de experiencia: {Título}{Empresa}{Ubicación}{Fechas}
\newcommand{\entrada}[4]{%
  \vspace{6pt}
  \begin{tabularx}{\textwidth}{@{}X r@{}}
    \textbf{#1} \textcolor{gris}{· #2} & \textcolor{grisCLaro}{\small\textit{#4}} \\
    \textcolor{grisCLaro}{\small\textit{#3}} & \\
  \end{tabularx}
  \vspace{2pt}
}

% Entrada de proyecto: {Nombre}{Tecnologías}
\newcommand{\proyecto}[2]{%
  \vspace{6pt}
  \textbf{#1} \quad \textcolor{gris}{\small\textit{| #2}}
  \vspace{2pt}
}

% Fila de habilidad: {Etiqueta}{Contenido}
\newcommand{\habilidad}[2]{%
  \textbf{#1:} #2 \\[2pt]
}

% ─────────────────────────────────────────────────────────────
\begin{document}
% ─────────────────────────────────────────────────────────────

% ── ENCABEZADO ────────────────────────────────────────────────
\begin{center}
  {\Huge\bfseries\color{azulcv} Jorge Pizarro Callejas} \\[6pt]
  {\large\color{gris} Ingeniero de Software \textbar{} Desarrollo Web \& Cloud} \\[6pt]
  {\small\color{gris}
    +56 976247629 \;·\;
    \href{mailto:jpizarrocallejas@gmail.com}{jpizarrocallejas@gmail.com} \;·\;
    \href{https://linkedin.com/in/jorgepizarrocallejas}{linkedin.com/in/jorgepizarrocallejas} \;·\;
    \href{https://github.com/jorgicio}{github.com/jorgicio} \;·\;
    Quilpué, Chile
  }
\end{center}

\vspace{4pt}

% ── PERFIL PROFESIONAL ────────────────────────────────────────
\section{Perfil Profesional}

Ingeniero de Ejecución en Informática con más de 10 años de experiencia en desarrollo de
software, especializado en aplicaciones web, servicios cloud (AWS) y automatización backend.
Experiencia comprobada en arquitecturas serverless con AWS Lambda, CI/CD con Jenkins, y
desarrollo fullstack con Java, Node.js, Angular y Vue.js. Orientado a la calidad del código,
trabajo en equipo remoto y resolución de problemas técnicos complejos.

% ── HABILIDADES TÉCNICAS ──────────────────────────────────────
\section{Habilidades Técnicas}

\habilidad{Lenguajes}{Java, JavaScript (Node.js), Python, PHP, C/C++, HTML5, CSS3, SQL}
\habilidad{Frameworks y librerías}{Angular, Vue.js, Django, Bootstrap, Tailwind CSS, Ionic, Expo Go}
\habilidad{Cloud y DevOps}{AWS Lambda, AWS CloudFormation, AWS CodeCommit, AWS CloudWatch, AWS DynamoDB, AWS Athena, Docker, Jenkins (CI/CD), SonarQube}
\habilidad{Bases de datos}{PostgreSQL, MySQL, AWS DynamoDB, AWS Athena}
\habilidad{Herramientas}{Git, GitHub, GitLab, VS Code, Linux, \LaTeX}
\habilidad{Metodologías}{CI/CD, control de versiones, QA automatizado, desarrollo remoto}

% ── EXPERIENCIA LABORAL ───────────────────────────────────────
\section{Experiencia Laboral}

\entrada{Ingeniero de Software Junior}{Tranciti (Waypoint Telecomunicaciones)}{Santiago, Chile (Remoto)}{Oct. 2021 -- Ene. 2025}
\begin{itemize}
  \item Desarrollo de microservicios y funciones serverless con AWS Lambda en Java y Node.js para una plataforma de geolocalización GPS.
  \item Implementación de infraestructura como código (IaC) con AWS CloudFormation y automatización de despliegues con Jenkins y AWS CodePipeline.
  \item Gestión de calidad de código con SonarQube e integración de Docker para entornos de desarrollo reproducibles.
  \item Desarrollo de módulo frontend para clientes en Angular, y soporte a apps móviles con Titanium, Ionic y Expo Go.
  \item Uso de PostgreSQL y AWS DynamoDB/Athena como motores de base de datos según el tipo de carga.
\end{itemize}

\entrada{Desarrollador Frontend}{Innovación y Proyectos Kimun Ltda.}{Temuco (Freelance Remoto)}{Dic. 2019}
\begin{itemize}
  \item Desarrollo y mejora del frontend de la plataforma web ``Mundo de vegetales'' utilizando HTML5, CSS3 y JavaScript.
\end{itemize}

\entrada{Desarrollador Web}{Pontificia Universidad Católica de Valparaíso}{Viña del Mar}{Sep. 2014 -- Abr. 2015}
\begin{itemize}
  \item Continuación del desarrollo de un sistema de aprendizaje de inglés para proyecto Conicyt usando PHP, jQuery y MySQL.
\end{itemize}

\entrada{Desarrollador Web}{Oxalix}{Viña del Mar}{Mar. 2013 -- Ene. 2014}
\begin{itemize}
  \item Desarrollo de formularios y módulos web con Python y Django como framework, integrando base de datos MySQL.
\end{itemize}

\entrada{Desarrollador Web y Encargado de Informática}{Del Puerto Comunicaciones}{Valparaíso}{Jun. 2011 -- Mar. 2014}
\begin{itemize}
  \item Desarrollo y mantenimiento de sitios web para clientes de agencia usando WordPress, HTML, CSS y jQuery. Práctica profesional de 360 horas continuada como empleado de planta.
\end{itemize}

\entrada{Desarrollador WordPress}{Serakin}{Viña del Mar (Freelance)}{Oct. 2016}
\begin{itemize}
  \item Desarrollo del sitio web corporativo del centro kinesiológico Serakin usando WordPress.
\end{itemize}

\entrada{Desarrollador Frontend}{Everis}{Santiago (Híbrido)}{Oct. -- Dic. 2015}
\begin{itemize}
  \item Desarrollo del panel de clientes del sitio web de Consorcio usando Bootstrap y jQuery.
\end{itemize}

% ── PROYECTOS PERSONALES ──────────────────────────────────────
\section{Proyectos Personales}

\proyecto{Bot de Indicadores Económicos (Telegram)}{Node.js, API REST}
\begin{itemize}
  \item Bot en Telegram que responde comandos para mostrar indicadores económicos en tiempo real usando la API de mindicador.cl.
\end{itemize}

\proyecto{Sitio web personal con portafolio}{Vue.js, Tailwind CSS, GraphQL}
\begin{itemize}
  \item Sitio web en desarrollo con carta de presentación, portafolio y blog. Frontend con Vue.js y Tailwind CSS, backend con GraphQL.
\end{itemize}

\proyecto{Tomython -- Temporizador Pomodoro}{Python, PyGTK, GStreamer}
\begin{itemize}
  \item Notificador Pomodoro de escritorio con notificaciones nativas de Linux, usando PyGTK y GStreamer.
\end{itemize}

\proyecto{Blog técnico en inglés}{Python, Pelican}
\begin{itemize}
  \item Blog estático creado con el generador Pelican en Python, publicado en inglés.
\end{itemize}

% ── EDUCACIÓN ─────────────────────────────────────────────────
\section{Educación}

\entrada{Ingeniería en Ejecución en Informática}{Universidad Técnica Federico Santa María}{Valparaíso, Chile}{2012 -- 2015}
\textit{\textcolor{gris}{Titulado.}}

\end{document}
